\chapter{Functors}

\begin{overviewbox}
A category is a world of objects connected by morphisms. A \emph{functor} is
a map between two such worlds that respects their structure. It sends objects
to objects and morphisms to morphisms, but not arbitrarily: it must preserve
composition and identity. A functor that ignores the structure --- sending
everything to a single point --- is still a functor; one that scrambles the
composition rules is not.

By the end of this chapter we will have seen that categories and functors
themselves form a category. This is one of the first signs that category theory
is ``self-referential'' in a productive way: the tools we build can be applied
to the framework that built them.
\end{overviewbox}


% ═══════════════════════════════════════════════════════════════════════════
\section{Structure-Preserving Maps}
\label{sec:structure-preserving}
% ═══════════════════════════════════════════════════════════════════════════

\subsection{Informally: Translating Between Categories}

Suppose you have two categories. A functor from the first to the second is a
systematic translation: every object in the first world is sent to some object
in the second, and every arrow in the first world is sent to an arrow in the
second --- in a way that respects the structure.

``Respects the structure'' has two precise requirements:
\begin{itemize}
  \item Composites translate to composites.
  \item Identity morphisms translate to identity morphisms.
\end{itemize}

If these two conditions hold, the translation is a functor.

\subsection{Expanding: A Functor from Recipes to Approvals}

We have seen the recipe category (Chop $\to$ Sauté $\to$ Add) and the
approval chain category (Request $\to$ Approved $\to$ Final). Both are
thin categories with three objects in a line. A functor between them is
a systematic renaming that preserves the ordering structure.

\begin{example}{The Status Functor}
Define $F$ as follows:

\medskip
\begin{center}
\begin{tabular}{ll}
  \textit{Recipe} & \textit{Approval chain} \\[4pt]
  \hline\\[-6pt]
  Chop & $\mapsto$ Request \\
  Sauté & $\mapsto$ Approved \\
  Add & $\mapsto$ Final \\
\end{tabular}
\end{center}
\medskip

On objects: $F(\mathrm{Chop}) = \mathrm{Request}$, and so on. On morphisms:
$F(f) = $ the unique morphism from Request to Approved (the manager's
approval). $F(g) = $ the unique morphism from Approved to Final. $F(g \circ f)
= $ the unique morphism from Request to Final.

Is this a functor? We check:
\begin{itemize}
  \item \textbf{Preserves composition}: $F(g \circ f) = F(g) \circ F(f)$.
        Both sides are the unique morphism from Request to Final. \checkmark
  \item \textbf{Preserves identity}: $F(\mathrm{id}_{\mathrm{Chop}}) =
        \mathrm{id}_{\mathrm{Request}}$. In a thin category, identity
        morphisms are the only morphisms from an object to itself, so this
        is automatic. \checkmark
\end{itemize}

$F$ is a functor.
\end{example}

\begin{example}{A Functor Between Orderings}
Consider the temperature ordering $\mathcal{T}$:
\[
  0°\text{C} \;\to\; 20°\text{C} \;\to\; 37°\text{C} \;\to\; 100°\text{C}
\]
and the size ordering $\mathcal{S}$: Small $\to$ Medium $\to$ Large.

Define $G : \mathcal{T} \to \mathcal{S}$ by:
\[
  G(0°\text{C}) = \mathrm{Small},\quad
  G(20°\text{C}) = \mathrm{Small},\quad
  G(37°\text{C}) = \mathrm{Medium},\quad
  G(100°\text{C}) = \mathrm{Large}.
\]

Two temperature values (freezing and room temperature) both map to Small.
On morphisms: if $A \leq B$ in $\mathcal{T}$, then $G(A) \leq G(B)$ in
$\mathcal{S}$ (check: $G(0°) = G(20°) = \mathrm{Small} \leq \mathrm{Small}$,
and so on). Composition and identity are preserved because all hom-sets have
at most one element.

$G$ is a functor. It ``collapses'' two temperature values into one size value
without breaking the ordering structure.
\end{example}

\begin{dialog}
When I want to check that $F$ is a functor from $\mathcal{C}$ to
$\mathcal{D}$, I run two checks:
\begin{enumerate}
  \item \textbf{Composition}: For every composable pair $f : A \to B$ and
        $g : B \to C$ in $\mathcal{C}$: is $F(g \circ f) = F(g) \circ F(f)$?
  \item \textbf{Identity}: For every object $A$ in $\mathcal{C}$: is
        $F(\mathrm{id}_A) = \mathrm{id}_{F(A)}$?
\end{enumerate}
If both hold for all morphisms and objects, $F$ is a functor. Note that the
composition on the left of each equation is in $\mathcal{C}$, and on the
right it is in $\mathcal{D}$.
\end{dialog}

\subsection{Formalizing: The Two Conditions}

A \term{functor} $F : \mathcal{C} \to \mathcal{D}$ consists of:
\begin{enumerate}
  \item An \term{object assignment}: for each $A \in \mathrm{Ob}(\mathcal{C})$,
        an object $F(A) \in \mathrm{Ob}(\mathcal{D})$.
  \item A \term{morphism assignment}: for each $f \in \mathrm{Hom}_\mathcal{C}(A,B)$,
        a morphism $F(f) \in \mathrm{Hom}_\mathcal{D}(F(A), F(B))$.
\end{enumerate}

Satisfying:
\begin{itemize}
  \item \textbf{Functoriality (composition)}: $F(g \circ f) = F(g) \circ F(f)$
        for all composable $f, g$ in $\mathcal{C}$.
  \item \textbf{Functoriality (identity)}: $F(\mathrm{id}_A) = \mathrm{id}_{F(A)}$
        for all $A \in \mathrm{Ob}(\mathcal{C})$.
\end{itemize}

We say $F$ is a functor \emph{from} $\mathcal{C}$ \emph{to} $\mathcal{D}$,
or a \term{covariant functor}, or just a functor. ($\mathcal{C}$ is the
\term{source} or \term{domain}; $\mathcal{D}$ is the \term{target} or
\term{codomain}.)

\subsection{Reorganizing: What the Conditions Enforce}

The two conditions enforce that a functor does not invent or destroy structure:

\textbf{Composition condition.} If, in $\mathcal{C}$, the path
$A \xrightarrow{f} B \xrightarrow{g} C$ exists and composes to $h : A \to C$,
then in $\mathcal{D}$ the path $F(A) \xrightarrow{F(f)} F(B) \xrightarrow{F(g)}
F(C)$ exists and composes to $F(h) = F(g) \circ F(f)$. Every composable
triangle in $\mathcal{C}$ maps to a composable triangle in $\mathcal{D}$,
and the triangles commute.

\textbf{Identity condition.} Every identity loop $A \xrightarrow{\mathrm{id}_A} A$
in $\mathcal{C}$ maps to the identity loop $F(A) \xrightarrow{\mathrm{id}_{F(A)}}
F(A)$ in $\mathcal{D}$. Functors do not turn identity morphisms into
non-trivial morphisms.

Together: \emph{a functor sends commutative diagrams to commutative diagrams}.

\subsection{Simplifying: The Functoriality Conditions}

\[
  F(g \circ f) = F(g) \circ F(f)
\]
\[
  F(\mathrm{id}_A) = \mathrm{id}_{F(A)}
\]

The first equation says composition commutes with $F$. The second says
identity commutes with $F$. These two equations are the complete definition
of a functor (given the object and morphism assignments).

\subsection{Formally: Functors and Commutative Diagrams}

\formalnote{Functors preserve commutativity}
If a diagram in $\mathcal{C}$ commutes, its image under any functor
$F : \mathcal{C} \to \mathcal{D}$ commutes in $\mathcal{D}$. This follows
directly from the functoriality conditions: every path in the original
diagram corresponds to an equal composite, and $F$ maps that equality to
an equality of composites in $\mathcal{D}$.

\formalnote{Functors and isomorphisms}
If $f : A \to B$ is an isomorphism in $\mathcal{C}$, then $F(f) : F(A) \to
F(B)$ is an isomorphism in $\mathcal{D}$, with inverse $F(f^{-1})$.
Proof: $F(f^{-1}) \circ F(f) = F(f^{-1} \circ f) = F(\mathrm{id}_A) =
\mathrm{id}_{F(A)}$, and similarly on the other side.

Functors \emph{preserve} isomorphisms. They may also \emph{create} them:
two non-isomorphic objects in $\mathcal{C}$ can have isomorphic images in
$\mathcal{D}$.


% ═══════════════════════════════════════════════════════════════════════════
\section{More Examples of Functors}
\label{sec:functor-examples}
% ═══════════════════════════════════════════════════════════════════════════

\subsection{Informally: Functors Are Everywhere}

Once you have the definition, you start seeing functors in many places. Some
translate faithfully, preserving all distinctions. Some collapse distinctions.
Some ignore most of the source and map everything to a single point. All of
these are functors, provided the two conditions hold.

\subsection{Expanding: Four Functors}

\begin{example}{The Identity Functor}
For any category $\mathcal{C}$, define $\mathrm{Id}_\mathcal{C} :
\mathcal{C} \to \mathcal{C}$ by:
\[
  \mathrm{Id}_\mathcal{C}(A) = A, \qquad \mathrm{Id}_\mathcal{C}(f) = f.
\]
Everything goes to itself. Both functoriality conditions hold trivially.
This is the ``do nothing'' functor on $\mathcal{C}$.
\end{example}

\begin{example}{A Constant Functor}
For any categories $\mathcal{C}$ and $\mathcal{D}$, and any fixed object
$D \in \mathcal{D}$, define $\Delta_D : \mathcal{C} \to \mathcal{D}$ by:
\[
  \Delta_D(A) = D \text{ for all } A, \qquad
  \Delta_D(f) = \mathrm{id}_D \text{ for all } f.
\]
Every object maps to $D$, every morphism maps to $\mathrm{id}_D$.

Check: $\Delta_D(g \circ f) = \mathrm{id}_D = \mathrm{id}_D \circ \mathrm{id}_D
= \Delta_D(g) \circ \Delta_D(f)$. \checkmark\;
$\Delta_D(\mathrm{id}_A) = \mathrm{id}_D = \mathrm{id}_{\Delta_D(A)}$. \checkmark

The constant functor collapses all of $\mathcal{C}$ to a single point in
$\mathcal{D}$. It is the most forgetful functor possible.
\end{example}

\begin{example}{A Functor from $\mathbf{1}$ to $\mathcal{C}$}
A functor $F : \mathbf{1} \to \mathcal{C}$ (from the one-object category
to $\mathcal{C}$) picks out one object in $\mathcal{C}$: whatever $F(*)$
is. The only morphism in $\mathbf{1}$ is $\mathrm{id}_*$, and it must map
to $\mathrm{id}_{F(*)}$.

This means: \emph{functors from $\mathbf{1}$ are in exact correspondence with
objects of $\mathcal{C}$}. The category $\mathbf{1}$ is the ``object selector.''
A functor from $\mathbf{1}$ is a way of pointing at a single object.
\end{example}

\begin{example}{A Functor from $\mathbf{2}$ to $\mathcal{C}$}
A functor $F : \mathbf{2} \to \mathcal{C}$ picks out a single morphism in
$\mathcal{C}$: the image of $\iota : 0 \to 1$. The two identity morphisms
of $\mathbf{2}$ must map to identities, and $F(\iota)$ can be any morphism
$F(0) \to F(1)$.

\emph{Functors from $\mathbf{2}$ are in exact correspondence with morphisms
of $\mathcal{C}$}. The category $\mathbf{2}$ is the ``morphism selector.''

This is why $\mathbf{2}$ is called the ``walking morphism'': it is the
minimal category containing a non-trivial morphism, and any morphism in any
category can be captured by a functor from it.
\end{example}

\subsection{Formalizing: Forgetful and Inclusion Functors}

Some functors appear so often they deserve names.

A \term{forgetful functor} is one that ``forgets'' some of the structure of
its source. We have not yet built enough structure to give a formal example,
but the constant functor $\Delta_D$ is an extreme case: it forgets everything.
The temperature-to-size functor $G$ from the previous section is a mild
forgetfulness: it forgets which of $0°\text{C}$ and $20°\text{C}$ you started
at.

An \term{inclusion functor} is one that embeds a subcategory into a larger
category without changing any object or morphism. For example, any category
$\mathcal{C}$ that contains the recipe category as a subcategory has an
inclusion functor $\mathrm{Recipe} \hookrightarrow \mathcal{C}$ that is the
identity on all objects and morphisms of $\mathrm{Recipe}$.

\subsection{Reorganizing: What Functors See and Ignore}

\begin{center}
\begin{tabular}{lll}
  Functor & What it sees & What it ignores \\[4pt]
  \hline\\[-6pt]
  $\mathrm{Id}_\mathcal{C}$ & everything & nothing \\
  $\Delta_D$ & nothing & everything \\
  Temperature $\to$ Size & ordering & which of two temps is cooler \\
  $\mathbf{1} \to \mathcal{C}$ & one object & all other objects and morphisms \\
  $\mathbf{2} \to \mathcal{C}$ & one morphism & all other morphisms \\
\end{tabular}
\end{center}

\medskip
Every functor lies on a spectrum from ``sees everything'' to ``sees nothing.''
The two conditions (preserve composition, preserve identity) constrain what it
can ignore, but leave a great deal of freedom.

\subsection{Simplifying: Functors as Translations}

A functor $F : \mathcal{C} \to \mathcal{D}$ is a translation of $\mathcal{C}$-
language into $\mathcal{D}$-language. The two conditions guarantee that every
sentence that is true in $\mathcal{C}$ (every commutative diagram) translates
to a true sentence in $\mathcal{D}$. The translation may collapse distinctions
but cannot introduce falsehoods.

\subsection{Formally: Faithful, Full, and Essentially Surjective}

\formalnote{Faithfulness and fullness}
A functor $F : \mathcal{C} \to \mathcal{D}$ is:
\begin{itemize}
  \item \term{faithful} if for all $A, B \in \mathcal{C}$, the function
        $F : \mathrm{Hom}_\mathcal{C}(A,B) \to \mathrm{Hom}_\mathcal{D}(F(A),F(B))$
        is injective (no two distinct morphisms get identified).
  \item \term{full} if that function is surjective (every morphism between
        the images is hit).
  \item \term{fully faithful} if it is both.
\end{itemize}

A faithful functor ``doesn't identify'' morphisms; a full one ``doesn't miss''
morphisms between images. The identity functor is fully faithful. The constant
functor is neither (if $\mathcal{C}$ has more than one morphism).

\formalnote{Essentially surjective}
$F$ is \term{essentially surjective} if for every $D \in \mathcal{D}$, there
exists $C \in \mathcal{C}$ with $F(C) \cong D$. Combined with full
faithfulness, this gives an \term{equivalence of categories} --- the
categorical substitute for isomorphism between categories, which we will
develop in a later chapter.


% ═══════════════════════════════════════════════════════════════════════════
\section{Composition of Functors}
\label{sec:functor-composition}
% ═══════════════════════════════════════════════════════════════════════════

\subsection{Informally: Translating Twice}

If $F$ translates from $\mathcal{C}$ to $\mathcal{D}$, and $G$ translates
from $\mathcal{D}$ to $\mathcal{E}$, then following $F$ then $G$ gives a
translation directly from $\mathcal{C}$ to $\mathcal{E}$. This composite
translation is a functor, called the \term{composite functor} $G \circ F$.

\subsection{Expanding: Composing the Examples}

We have $F : \mathrm{Recipe} \to \mathrm{Approval}$ (the status functor from
§3.1) and suppose $H : \mathrm{Approval} \to \mathcal{S}$ maps each approval
stage to a project phase (Request $\mapsto$ Planning, Approved $\mapsto$
Building, Final $\mapsto$ Done).

The composite $H \circ F : \mathrm{Recipe} \to \mathcal{S}$ maps:
\[
  \mathrm{Chop} \mapsto \mathrm{Planning},\quad
  \mathrm{Sauté} \mapsto \mathrm{Building},\quad
  \mathrm{Add} \mapsto \mathrm{Done}.
\]

\begin{dialog}
I check that $H \circ F$ is a functor:
\begin{itemize}
  \item $(H \circ F)(g \circ f) = H(F(g \circ f)) = H(F(g) \circ F(f)) =
        H(F(g)) \circ H(F(f)) = (H \circ F)(g) \circ (H \circ F)(f)$. \checkmark
  \item $(H \circ F)(\mathrm{id}_A) = H(F(\mathrm{id}_A)) = H(\mathrm{id}_{F(A)})
        = \mathrm{id}_{H(F(A))} = \mathrm{id}_{(H \circ F)(A)}$. \checkmark
\end{itemize}
Each step uses the functoriality of $F$ (first) and $G$ (second). The
argument works for any two composable functors.
\end{dialog}

\subsection{Formalizing: The Composite Functor}

Given functors $F : \mathcal{C} \to \mathcal{D}$ and $G : \mathcal{D} \to
\mathcal{E}$, their \term{composite} $G \circ F : \mathcal{C} \to \mathcal{E}$
is defined by:
\[
  (G \circ F)(A) = G(F(A)),\qquad (G \circ F)(f) = G(F(f)).
\]

\textbf{$G \circ F$ is a functor.} Functoriality for composition:
\[
  (G \circ F)(g \circ f) = G(F(g \circ f)) = G(F(g) \circ F(f))
  = G(F(g)) \circ G(F(f)) = (G \circ F)(g) \circ (G \circ F)(f).
\]
Functoriality for identity: $(G \circ F)(\mathrm{id}_A) = G(F(\mathrm{id}_A))
= G(\mathrm{id}_{F(A)}) = \mathrm{id}_{G(F(A))} = \mathrm{id}_{(G \circ F)(A)}$.

\textbf{Composition of functors is associative.} For $F : \mathcal{C} \to
\mathcal{D}$, $G : \mathcal{D} \to \mathcal{E}$, $H : \mathcal{E} \to
\mathcal{F}$:
\[
  (H \circ G) \circ F = H \circ (G \circ F).
\]
This holds because it holds pointwise on objects and morphisms (function
composition is associative).

\textbf{Identity functors are neutral.} $\mathrm{Id}_\mathcal{D} \circ F =
F = F \circ \mathrm{Id}_\mathcal{C}$.

\subsection{Reorganizing: Functors Are Morphisms}

We now have:
\begin{itemize}
  \item \textbf{Objects}: categories $\mathcal{C}$, $\mathcal{D}$, $\ldots$
  \item \textbf{Morphisms}: functors $F : \mathcal{C} \to \mathcal{D}$
  \item \textbf{Composition}: $G \circ F : \mathcal{C} \to \mathcal{E}$,
        associative
  \item \textbf{Identity}: $\mathrm{Id}_\mathcal{C} : \mathcal{C} \to
        \mathcal{C}$, neutral
\end{itemize}

This matches the pattern for a category exactly. Categories and functors
form a category.

\subsection{Simplifying: Functor Composition}

\[
  (G \circ F)(A) = G(F(A)),\qquad (G \circ F)(f) = G(F(f))
\]
\[
  (H \circ G) \circ F = H \circ (G \circ F),\qquad
  \mathrm{Id}_\mathcal{D} \circ F = F = F \circ \mathrm{Id}_\mathcal{C}
\]

These are the same associativity and unit equations as for morphisms in any
category --- because functors \emph{are} morphisms, in the category of
categories.

\subsection{Formally: The Category $\mathbf{Cat}$}

\formalnote{Definition}
The \term{category of categories} $\mathbf{Cat}$ has:
\begin{itemize}
  \item Objects: (small) categories.
  \item Morphisms: functors between them.
  \item Composition: composite functor $G \circ F$.
  \item Identity: identity functor $\mathrm{Id}_\mathcal{C}$.
\end{itemize}
The axioms are verified above. $\mathbf{Cat}$ is a category.

\formalnote{Size remark}
To avoid set-theoretic paradoxes (a ``category of all categories'' including
itself), $\mathbf{Cat}$ is typically defined to have \emph{small} categories
as objects. The category of all locally small categories is sometimes called
$\mathbf{CAT}$. We will not belabor this point, but the reader should be
aware that size considerations are lurking.

\formalnote{Endofunctors}
A functor $F : \mathcal{C} \to \mathcal{C}$ from a category to itself is
called an \term{endofunctor}. Endofunctors can be composed repeatedly:
$F^n = F \circ F \circ \cdots \circ F$ ($n$ times). Endofunctors on a
category $\mathcal{C}$ form a category $[\mathcal{C}, \mathcal{C}]$ (a
\term{functor category}, to be developed in Chapter~4).


% ═══════════════════════════════════════════════════════════════════════════
\section{Contravariant Functors}
\label{sec:contravariant}
% ═══════════════════════════════════════════════════════════════════════════

\subsection{Informally: Going Backwards}

Some translations between categories send morphisms in the \emph{opposite}
direction. An arrow $f : A \to B$ in the source gets sent to an arrow
$F(f) : F(B) \to F(A)$ in the target --- the direction is reversed.

Such a translation still has to respect composition and identity, but the
composition rule looks like a reversal. This is a \term{contravariant functor}.

\subsection{Expanding: The Reversal Pattern}

\begin{example}{An Opposite Ordering}
Consider the temperature ordering $\mathcal{T}$ (colder to warmer) and the
size ordering $\mathcal{S}^{\mathrm{op}}$ (larger to smaller: Large $\to$
Medium $\to$ Small).

Define $F : \mathcal{T} \to \mathcal{S}^{\mathrm{op}}$ by: colder things
are \emph{larger} in $\mathcal{S}^{\mathrm{op}}$:
\[
  F(0°\text{C}) = \mathrm{Large},\quad F(20°\text{C}) = \mathrm{Medium},
  \quad F(37°\text{C}) = \mathrm{Small},\quad F(100°\text{C}) = \mathrm{Small}.
\]
A morphism $0°\text{C} \to 37°\text{C}$ in $\mathcal{T}$ (colder to warmer)
maps to Large $\to$ Small in $\mathcal{S}^{\mathrm{op}}$ (larger to smaller).
The direction of the ordering is preserved in the opposite sense.

This is a covariant functor $\mathcal{T} \to \mathcal{S}^{\mathrm{op}}$, or
equivalently a contravariant functor $\mathcal{T} \to \mathcal{S}$.
\end{example}

\begin{dialog}
I see a map $F$ that sends morphisms ``backwards'': $f : A \to B$ maps to
$F(f) : F(B) \to F(A)$. I ask:

\begin{itemize}
  \item Does $F(g \circ f) = F(f) \circ F(g)$? (Note the reversed order.)
  \item Does $F(\mathrm{id}_A) = \mathrm{id}_{F(A)}$?
\end{itemize}

If yes to both, $F$ is a contravariant functor. The reversal in composition
is not a mistake --- it is forced by the direction flip.
\end{dialog}

\subsection{Formalizing: Contravariance as a Functor from the Opposite}

A \term{contravariant functor} $F : \mathcal{C} \to \mathcal{D}$ is exactly
a (covariant) functor $F : \mathcal{C}^{\mathrm{op}} \to \mathcal{D}$.

On objects: $F$ assigns an object $F(A) \in \mathcal{D}$ to each $A \in
\mathcal{C}$. On morphisms: for $f : A \to B$ in $\mathcal{C}$, $F$ assigns
$F(f) : F(B) \to F(A)$ in $\mathcal{D}$.

Functoriality:
\[
  F(g \circ f) = F(f) \circ F(g) \qquad\text{(note reversed order)}
\]
\[
  F(\mathrm{id}_A) = \mathrm{id}_{F(A)}
\]

This is the covariant functoriality applied to $\mathcal{C}^{\mathrm{op}}$:
in $\mathcal{C}^{\mathrm{op}}$, composition of $f^{\mathrm{op}}$ and
$g^{\mathrm{op}}$ gives $(f \circ g)^{\mathrm{op}}$, and $F$ applied to
$(f \circ g)^{\mathrm{op}}$ gives $F(f \circ g) = F(g) \circ F(f)$... wait,
applying the covariant rule to $\mathcal{C}^{\mathrm{op}}$ correctly gives
$F(f^{\mathrm{op}} \circ_{\mathrm{op}} g^{\mathrm{op}}) = F(f^{\mathrm{op}})
\circ F(g^{\mathrm{op}})$, which translates back to $F(g \circ f) =
F(f) \circ F(g)$ in $\mathcal{D}$.

\subsection{Reorganizing: Contravariance = Covariance on the Opposite}

We do not need a new concept for contravariant functors. Every contravariant
functor $\mathcal{C} \to \mathcal{D}$ is a covariant functor
$\mathcal{C}^{\mathrm{op}} \to \mathcal{D}$. The opposite category absorbs
the reversal.

This is the first example of the \emph{duality principle} in practice: instead
of developing a parallel theory for ``backwards'' maps, we build the reversal
into the objects (via $\mathcal{C}^{\mathrm{op}}$) and reuse the existing
theory.

\subsection{Simplifying: Contravariance in One Line}

A contravariant functor from $\mathcal{C}$ to $\mathcal{D}$ is a covariant
functor:
\[
  F : \mathcal{C}^{\mathrm{op}} \longrightarrow \mathcal{D}.
\]

\subsection{Formally: Bifunctors}

\formalnote{Bifunctors}
A \term{bifunctor} is a functor of two arguments: $F : \mathcal{C} \times
\mathcal{D} \to \mathcal{E}$, where $\mathcal{C} \times \mathcal{D}$ is
the \term{product category} whose objects are pairs $(C, D)$ and whose
morphisms are pairs $(f, g)$. A mixed-variance bifunctor has the form
$F : \mathcal{C}^{\mathrm{op}} \times \mathcal{D} \to \mathcal{E}$:
contravariant in the first argument, covariant in the second. The
hom-functor $\mathrm{Hom}(-, -)$ (to be studied in Chapter~5) is the
canonical example.


% ═══════════════════════════════════════════════════════════════════════════
\section*{Chapter Summary}
\addcontentsline{toc}{section}{Chapter Summary}
% ═══════════════════════════════════════════════════════════════════════════

\begin{summarybox}
A \textbf{functor} $F : \mathcal{C} \to \mathcal{D}$ assigns objects to
objects and morphisms to morphisms, satisfying:
\[
  F(g \circ f) = F(g) \circ F(f), \qquad F(\mathrm{id}_A) = \mathrm{id}_{F(A)}.
\]

Functors \textbf{preserve commutativity} of diagrams and \textbf{preserve
isomorphisms}.

\textbf{Composition of functors} $(G \circ F)(A) = G(F(A))$ is associative
and has the identity functor as neutral element.

\textbf{Categories and functors form a category}, called $\mathbf{Cat}$.

A \textbf{contravariant functor} $\mathcal{C} \to \mathcal{D}$ is a covariant
functor $\mathcal{C}^{\mathrm{op}} \to \mathcal{D}$.

\textbf{Special functors:} identity ($\mathrm{Id}_\mathcal{C}$), constant
($\Delta_D$), faithful (injective on hom-sets), full (surjective on hom-sets).
Functors from $\mathbf{1}$ pick out objects; functors from $\mathbf{2}$ pick
out morphisms.
\end{summarybox}

% ═══════════════════════════════════════════════════════════════════════════
\section*{Formal Restatement}
\addcontentsline{toc}{section}{Formal Restatement}
% ═══════════════════════════════════════════════════════════════════════════

\begin{formalrestatement}
\textbf{Functor.}\quad $F : \mathcal{C} \to \mathcal{D}$ assigns
$F : \mathrm{Ob}(\mathcal{C}) \to \mathrm{Ob}(\mathcal{D})$ and
$F : \mathrm{Hom}_\mathcal{C}(A,B) \to \mathrm{Hom}_\mathcal{D}(F(A),F(B))$,
satisfying:
\[
  F(g \circ f) = F(g) \circ F(f), \qquad F(\mathrm{id}_A) = \mathrm{id}_{F(A)}.
\]

\medskip
\textbf{Composition.}\quad $(G \circ F)(A) = G(F(A))$, $(G \circ F)(f) =
G(F(f))$; associative; $\mathrm{Id}_\mathcal{D} \circ F = F =
F \circ \mathrm{Id}_\mathcal{C}$.

\medskip
\textbf{Properties.}\quad $F$ is \emph{faithful} iff injective on each
hom-set; \emph{full} iff surjective on each hom-set; \emph{essentially
surjective} iff every $D \in \mathcal{D}$ satisfies $D \cong F(C)$ for some
$C$.

\medskip
\textbf{Contravariant.}\quad A contravariant functor $\mathcal{C} \to
\mathcal{D}$ is a covariant functor $\mathcal{C}^{\mathrm{op}} \to
\mathcal{D}$; sends $f : A \to B$ to $F(f) : F(B) \to F(A)$; satisfies
$F(g \circ f) = F(f) \circ F(g)$.

\medskip
\textbf{$\mathbf{Cat}$.}\quad The category of small categories and functors:
objects are small categories, morphisms are functors, composition and identity
as above.

\medskip
\noindent\textit{Chapter~4 introduces natural transformations: morphisms
between functors, which complete the three-level hierarchy of category theory.}
\end{formalrestatement}
